% ===================================================================
\chapter{Interpreting the Hardware Numbers}
\label{ch:hwinterpret}

Kiem reports a comparable set of hardware numbers for his DRHE compressor
\cite{kiem2025}, and set side by side the two sets differ on almost every line.
A difference on its own explains nothing, so this chapter puts the numbers on
the same footing in three steps. Every design is implemented through placement
and routing, because Kiem's figures are post-route
(Section~\ref{sec:hw_counting}). Kiem's compressor is rebuilt in this thesis's
tool flow (Section~\ref{sec:hw_replica}). And the path from this thesis's
datapath to the rebuilt design is taken one design decision at a time, with
every step implemented (Section~\ref{sec:hw_ladder}). The remaining sections
read latency, throughput and resources against the sensor and the device, and
the chapter closes with a ledger of every difference from Kiem
(Section~\ref{sec:hw_ledger}).

Kiem's sensor has one transmitter, so his lag-1 predictor is the
per-transmitter predictor for that sensor. The comparison therefore starts
from \emph{DRHE-1}: this thesis's DRHE datapath predicting at lag~1. Its
datapath is that of the final lag-$\nTx$ compressor; only the indexing of the
predictor state differs.

% -------------------------------------------------------------------
\section{Estimates against routed results}
\label{sec:hw_counting}

HLS reports what it \emph{expects} a circuit to cost; Vivado reports what the
placed and routed circuit \emph{does} cost. Each core was implemented on its
own (out of context) on the \texttt{xcku5p-ffvb676-2-e} at
\SI{10}{\nano\second}, with \texttt{vivado.flow=impl} in its HLS
configuration. Out of context and in the board system, the DRHE compressor's
LUTs and flip-flops agree to within \SI{2}{\percent} and its DSPs and block RAM
exactly; LPC's block RAM differs by about \SI{10}{\percent}, because Vivado maps
its memories differently. The in-system figures of Table~\ref{tab:resources}
are the ones quoted for the final cores. Table~\ref{tab:hw_hlsvsimpl} compares
the two kinds of number.

\begin{table}[htbp]
\centering
\caption[HLS estimates against implemented results]{HLS estimates against the
same compressors after placement and routing, out of context, on the
\texttt{xcku5p} with a \SI{10}{\nano\second} constraint. BRAM in 36\,Kb blocks.
$F_{\max} = 1/(10\,\mathrm{ns} - \mathrm{WNS})$ after routing.}
\label{tab:hw_hlsvsimpl}
\small
\setlength{\tabcolsep}{3.2pt}
\begin{tabular}{@{}l rr rr rr rr rr@{}}
\toprule
& \multicolumn{2}{c}{LUT} & \multicolumn{2}{c}{FF} & \multicolumn{2}{c}{DSP} & \multicolumn{2}{c}{BRAM} & \multicolumn{2}{c}{$F_{\max}$ (MHz)}\\
\cmidrule(lr){2-3}\cmidrule(lr){4-5}\cmidrule(lr){6-7}\cmidrule(lr){8-9}\cmidrule(lr){10-11}
Design & HLS & impl. & HLS & impl. & HLS & impl. & HLS & impl. & HLS & impl.\\
\midrule
DRHE, lag $\nTx$ & \HW{drhen}{hLUT} & \textbf{\HW{drhen}{LUT}} & \HW{drhen}{hFF} & \textbf{\HW{drhen}{FF}} & \HW{drhen}{hDSP} & \textbf{\HW{drhen}{DSP}} & 40 & \textbf{\HW{drhen}{BRAMT}} & 81.9 & \textbf{\HW{drhen}{fmax}}\\
LPC, lags $\nTx, 2\nTx$ & \HW{lpcn}{hLUT} & \textbf{\HW{lpcn}{LUT}} & \HW{lpcn}{hFF} & \textbf{\HW{lpcn}{FF}} & \HW{lpcn}{hDSP} & \textbf{\HW{lpcn}{DSP}} & 96 & \textbf{\HW{lpcn}{BRAMT}} & 85.4 & \textbf{\HW{lpcn}{fmax}}\\
DRHE-1 (lag 1)   & \HW{drheone}{hLUT} & \HW{drheone}{LUT} & \HW{drheone}{hFF} & \HW{drheone}{FF} & \HW{drheone}{hDSP} & \HW{drheone}{DSP} & 20 & \HW{drheone}{BRAMT} & 81.9 & \HW{drheone}{fmax}\\
\bottomrule
\end{tabular}
\end{table}

HLS over-estimates logic by about three times (\HW{drhen}{estLUT} for DRHE,
\HW{lpcn}{estLUT} for LPC). Most of the excess is logic HLS generates itself
from C++. It costs each of the packer's eight inserts into the 512-bit register
as a full 512-bit barrel shifter, although the value shifted is at most 29 bits
wide and Vivado removes most of each shifter; and it costs each floating-point
\texttt{atan2}, a CORDIC built from generic C++, at \num{20129} LUTs against
about \num{5980} implemented. DSPs and block RAM, which are fixed blocks HLS
allocates itself, are estimated almost exactly. The clock estimate is
pessimistic in the same way: HLS predicts \SI{81.9}{\mega\hertz} and
\SI{85.4}{\mega\hertz}, yet every design meets \SI{100}{\mega\hertz} after
routing, and the board systems do so with \SIrange{0.17}{0.34}{\nano\second}
of slack (Table~\ref{tab:builds}). HLS estimates are useful for comparing
versions of one design; every number in this chapter compared with Kiem's is
post-route.

% -------------------------------------------------------------------
\section{Rebuilding the reference design}
\label{sec:hw_replica}

This thesis's DRHE and Kiem's differ in at least seven ways at once:
prediction axis, arithmetic, bit packer, state reset, feedback structure, tool
version (Vitis HLS 2026.1 against 2022.2) and device. To separate them, his
compressor was rebuilt from the description in his Section~4.3.4: his
fixed-point types (magnitude \texttt{ap\_ufixed<16,16>}, phase
\texttt{ap\_fixed<16,3>}, intermediates \texttt{ap\_fixed<32,4>}), the
\texttt{hls::sqrt}, \texttt{atan2}, \texttt{sin} and \texttt{cos} he names,
twenty $1024$-word state arrays driven from the input sample (open loop), a
58-bit word per channel packed into the 256-bit output, and all four channels
at $\mathrm{II} = 1$. He does not describe how the state is cleared at the
start of a frame; the replica reads it as zero on the first ramp, which is the
only reading consistent with $\mathrm{II} = 1$ (Section~\ref{sec:hw_ladder}).
The replica and DRHE-1 are one source file with four compile-time switches
(packer, state reset, arithmetic, loop), and every variant decodes ColoRadar
losslessly in C-simulation.

\begin{table}[htbp]
\centering
\caption[The replica against Kiem's published design]{The rebuilt design
against Kiem's published compression IP (his Table~5.3 \cite{kiem2025}).
Replica figures are post-route, out of context, on the \texttt{xcku5p} at
speed grades $-2$ and $-1$; Kiem's are post-route, in system, on the
\texttt{XCZU3CG} at speed grade $-1$. BRAM in 36\,Kb blocks.}
\label{tab:hw_replica}
\small
\begin{tabular}{@{}lrrr@{}}
\toprule
& Replica ($-2$) & Replica ($-1$) & Kiem \cite{kiem2025}\\
\midrule
LUT & \HW{ablk}{LUT} & \HW{ablksgone}{LUT} & \num{45892}\\
FF & \HW{ablk}{FF} & \HW{ablksgone}{FF} & \num{9243}\\
DSP & \HW{ablk}{DSP} & \HW{ablksgone}{DSP} & \num{44}\\
BRAM & \HW{ablk}{BRAMT} & \HW{ablksgone}{BRAMT} & \num{10}\\
Initiation interval & \HW{ablk}{hII} & \HW{ablksgone}{hII} & 1\\
Pipeline depth (HLS, cycles) & \HW{ablk}{hDepth} & \HW{ablksgone}{hDepth} & 23\\
Achieved clock (MHz) & \HW{ablk}{fmax} & \HW{ablksgone}{fmax} & 100 (met)\\
\bottomrule
\end{tabular}
\end{table}

The replica matches Kiem where his design is fully specified: block RAM
exactly, DSPs closely, $\mathrm{II} = 1$, and \SI{100}{\mega\hertz} on a
speed-grade $-1$ part (Table~\ref{tab:hw_replica}). It does not reproduce his
LUT and flip-flop counts, which his design as described does not need. Three
explanations were tested by building each into the replica: his slower speed
grade, the most literal reading of his one-sentence packer description, and
his deeper pipeline. If their effects add, they account for about
\SI{45}{\percent} of the flip-flop gap and \SI{5}{\percent} of the LUT gap.
The rest most likely lies in the four-year-older fixed-point math library of
his tool version, or in logic his thesis does not describe, such as the
debug probes his Figure~5.4 shows. His LUT figure therefore measures his
build, not the cost of the DRHE algorithm; the replica is the fair yardstick
for what this thesis's design decisions cost.

% -------------------------------------------------------------------
\section{From this thesis's datapath to Kiem's, one change at a time}
\label{sec:hw_ladder}

Table~\ref{tab:hw_ladder} is the central result of this chapter. Each row is a
complete lossless compressor implemented through routing; each changes one
design decision from the row above, from DRHE-1 to the replica. The difference
between adjacent rows is the cost of that one decision.

\begin{table}[htbp]
\centering
\caption[From DRHE-1 to the replica, one change per step]{From DRHE-1 to the
replica of Kiem's design, one design decision per step. Every row is lossless
and implemented through routing, out of context, on the \texttt{xcku5p} at
\SI{10}{\nano\second}. II and depth are from HLS. BRAM in 36\,Kb blocks.}
\label{tab:hw_ladder}
\small
\setlength{\tabcolsep}{4pt}
\begin{tabular}{@{}cl rrrr cc r@{}}
\toprule
Step & Change from the row above & LUT & FF & DSP & BRAM & II & Depth & MHz\\
\midrule
0 & DRHE-1 (float, closed loop, lag 1) & \HW{drheone}{LUT} & \HW{drheone}{FF} & \HW{drheone}{DSP} & \HW{drheone}{BRAMT} & \HW{drheone}{hII} & \HW{drheone}{hDepth} & \HW{drheone}{fmax}\\
1 & packer shift offset made unsigned & \HW{ablpone}{LUT} & \HW{ablpone}{FF} & \HW{ablpone}{DSP} & \HW{ablpone}{BRAMT} & \HW{ablpone}{hII} & \HW{ablpone}{hDepth} & \HW{ablpone}{fmax}\\
2 & Kiem's 58-bit per-channel packer & \HW{ablptwo}{LUT} & \HW{ablptwo}{FF} & \HW{ablptwo}{DSP} & \HW{ablptwo}{BRAMT} & \HW{ablptwo}{hII} & \HW{ablptwo}{hDepth} & \HW{ablptwo}{fmax}\\
3 & state cleared by read, not write & \HW{ablptwor}{LUT} & \HW{ablptwor}{FF} & \HW{ablptwor}{DSP} & \HW{ablptwor}{BRAMT} & \HW{ablptwor}{hII} & \HW{ablptwor}{hDepth} & \HW{ablptwor}{fmax}\\
4 & Kiem's fixed-point arithmetic & \HW{ablfxc}{LUT} & \HW{ablfxc}{FF} & \HW{ablfxc}{DSP} & \HW{ablfxc}{BRAMT} & \HW{ablfxc}{hII} & \HW{ablfxc}{hDepth} & \HW{ablfxc}{fmax}\\
5 & open loop $=$ replica of Kiem & \HW{ablk}{LUT} & \HW{ablk}{FF} & \HW{ablk}{DSP} & \HW{ablk}{BRAMT} & \HW{ablk}{hII} & \HW{ablk}{hDepth} & \HW{ablk}{fmax}\\
\midrule
--- & Kiem, published \cite{kiem2025} & \num{45892} & \num{9243} & \num{44} & \num{10} & 1 & 23 & 100\\
--- & DRHE, lag $\nTx$ (final design) & \HW{drhen}{LUT} & \HW{drhen}{FF} & \HW{drhen}{DSP} & \HW{drhen}{BRAMT} & \HW{drhen}{hII} & \HW{drhen}{hDepth} & \HW{drhen}{fmax}\\
\bottomrule
\end{tabular}
\end{table}

\paragraph{Block RAM is the word width.} DRHE's state is five values per range
bin and channel, twenty arrays. DRHE-1 stores them as 32-bit floats, so each
$1024$-word array fills a 36\,Kb block: 20 blocks. Kiem's 16-bit words fit two
arrays to a block: 10 blocks, which the replica matches exactly. The final
compressor keeps one word per transmitter and bin, $12 \times 128 = 1536$ words
per array, which HLS rounds up to two blocks: 40 in total. Against DRHE-1 it
changes nothing else --- the same 132 DSPs and the same schedule --- so at
ColoRadar's size the per-transmitter prediction costs 20 blocks of block RAM
and nothing more. The state grows with bins times transmitters: at Kiem's 512
bins and 12 transmitters it would be 160 blocks in floating point and 80 in his
16-bit types.

\paragraph{DSPs are floating point.} Every DSP in DRHE belongs to a
floating-point operator: 16 multipliers at 3 slices, 12 adders at 2, four
sine/cosine units at 13 and two arctangents at 4, 132 in all. A 16-bit
fixed-point product needs one slice and an addition none. Four channels would
need eight sine/cosine units; at $\mathrm{II} = 2$ HLS shares each unit
between two channels, and step~3 shows that without the sharing the float
design needs \HW{ablptwor}{DSP}.

\paragraph{The pipeline depth is the closed loop.} DRHE's pipeline is 58
cycles deep against 23 in Kiem's HLS report. It is two chains in series: the
prediction (state, filters, sine and cosine, rounding), then the polar
conversion that updates the state, whose \texttt{atan2} alone is 28 stages.
The second chain must wait for the first because DRHE updates its state from
the \emph{reconstructed} sample, which is what guarantees the encoder and
decoder never disagree (Section~\ref{sec:robust}). Driving the state from the
input instead lets the chains overlap and cuts the float depth to 32; fixed
point shortens both chains further, to \HW{ablfxc}{hDepth} closed loop and
\HW{ablk}{hDepth} open loop. Of DRHE's 58 cycles, 26 are the closed loop; at
\SI{260}{\nano\second} it is a cost this application does not notice.

\paragraph{The initiation interval is one write.} HLS reports why DRHE runs at
$\mathrm{II} = 2$: ``unable to schedule load operation \ldots due to limited
memory ports''. Each state array is a block RAM with two ports, and each cycle
needs one read and one write. DRHE also clears the state at the start of a
frame with an \emph{extra} write; it happens on only the first ramps, but HLS
schedules for the worst case, and three accesses on two ports take two cycles.
Step~3 clears the state by ignoring the read on the first ramp instead, and the
same float design schedules at $\mathrm{II} = 1$ --- at the price of losing the
operator sharing, which doubles its DSPs and grows its logic. It is the only
step that makes the design substantially larger, and it is why the final
design keeps $\mathrm{II} = 2$: with 24 times more throughput than four
ColoRadar channels need (Section~\ref{sec:hw_throughput}), half the rate buys
half the arithmetic.

\paragraph{LUTs and flip-flops.} Against Kiem's published figures, the final
compressor uses fewer LUTs (0.77$\times$) but more flip-flops (2.55$\times$).
The apparent contradiction comes from the unexplained gap between his build and
the replica, which is large for LUTs and small for flip-flops. Against the
replica, DRHE-1 is larger in both, by amounts the ladder accounts for
completely. The floating-point operators are \HW{drheone}{fopLUTpct}\,\% of
DRHE-1's LUTs, the two arctangents alone a third. Steps 1 and 2 show that
DRHE's packer costs about \num{4200} LUTs more than Kiem's: a signed shift
offset keeps logic for negative shifts that never happen, and eight inserts
into a 512-bit register cost more than one 232-bit insert. Step~4, fixed
point, removes most of the rest, and flip-flops fall further than LUTs because
the values in flight become both narrower and fewer. Step~5, the open loop, is
almost free in fixed point. Changing \emph{only} DRHE-1's arithmetic to Kiem's
types, keeping everything else, gives \HW{ablf}{LUT} LUTs (\HW{ablf}{pDLUT}\,\%
fewer), exactly Kiem's \HW{ablf}{DSP} DSPs and half the block RAM, and
compresses ColoRadar to within \SI{0.07}{\percent} of the floating-point design.
It meets \SI{100}{\mega\hertz} with little margin (\HW{ablf}{fmax}\,MHz),
because its short closed loop puts 26 carry-chain levels before the state
write; one or two more pipeline stages would fix that.

\paragraph{The clock.} Every design meets \SI{100}{\mega\hertz}. The replica
closes at \HW{ablksgone}{fmax}\,MHz on a speed-grade $-1$ part and
\HW{ablk}{fmax}\,MHz on a $-2$, and DRHE-1 closes at
\HW{ablctrlsgone}{fmax}\,MHz on a $-1$ part, so this thesis's floating-point
datapath meets Kiem's clock on Kiem's grade of silicon.

% -------------------------------------------------------------------
\section{Latency}
\label{sec:hw_latency}

Three quantities are called latency, and they answer different questions.

\paragraph{Pipeline latency} is how long one sample takes to pass through:
58 cycles, \SI{580}{\nano\second}, for DRHE, against Kiem's HLS estimate of
23 cycles. The difference is the closed loop and the latency of floating-point
operators (Section~\ref{sec:hw_ladder}). It is \SI{0.7}{\percent} of a
ColoRadar ramp interval and has no practical consequence here.

\paragraph{Frame time} is how long a whole frame takes when it is fed as fast
as the core accepts it: \SI{0.61}{\milli\second} for DRHE and
\SI{3.25}{\milli\second} for LPC on the board. It is really a throughput
figure. DRHE's is set by its initiation interval --- \num{24576} beats at two
cycles each, plus the pipeline draining once per ramp --- and LPC's by its
three passes (Section~\ref{sec:lpcimpl}). The sensor delivers a frame over
\SI{15.07}{\milli\second}, so DRHE is busy about \SI{4}{\percent} of the time.

\paragraph{Added latency} is how much later the downstream processor receives
the last bit of a frame because of compression, measured with Kiem's definition
from the last input word read to the last compressed word written. This is the
latency that matters. DRHE streams, so on the board only the data still in
flight remains: \SI{13.1}{\micro\second}. In a sensor-paced system no backlog
forms, and only the last ramp's 320 cycles, \SI{3.2}{\micro\second}, remain.
LPC cannot emit anything until it has seen the whole frame, so after the last
input word it still has its solve and emit passes to run:
\SI{3.0}{\milli\second}, however the frame is fed. That is a factor of 230 on
the board and about a thousand in a sensor-paced system. Kiem's
\SI{79.4}{\micro\second} uses the same definition but includes the range FFT
of the last ramp and its DMA transfers. Measured over the same chain, with the
range FFT on chip (Section~\ref{sec:feresults}), this thesis's DRHE takes
\SI{21.3}{\micro\second}; both are tens of microseconds and negligible against
the frame. The
decoders' latencies, \SI{339}{\micro\second} and \lpcDecLat, are their own
backlog when the compressed input ends, not a dependence on the whole frame.

\paragraph{When latency becomes a problem.} There are three thresholds. The
hardest is throughput: a compressor slower than the sensor builds a backlog
and loses data. DRHE uses 320 of the \num{7850} cycles each ColoRadar ramp
provides. The second is the system's schedule. Kiem allotted
\SI{15}{\milli\second} of his \SI{50}{\milli\second} cycle to compression; on
ColoRadar, at \SI{5}{\hertz}, both designs use a tiny fraction of the gap after
acquisition. Scaled to Kiem's frame, 21 times larger, DRHE still keeps up ramp
by ramp (\SI{10.9}{\micro\second} of work per \SI{20.5}{\micro\second} ramp),
but LPC's post-ingest passes grow to about \SI{64}{\milli\second}, four times
his budget and longer than his whole cycle, and its frame buffer to
\SI{67}{\mega bit}, more than the block RAM of either device. Even the
one-third saving available from overlapping its passes would leave it near
\SI{45}{\milli\second}. The third threshold is the application: at
\SI{100}{\kilo\metre\per\hour} LPC's \SI{3.0}{\milli\second} is
\SI{8.4}{\centi\metre} of travel and DRHE's \SI{13.1}{\micro\second} a third of
a millimetre, both small next to a \SI{5}{\hertz} radar's own
\SI{200}{\milli\second}. On ColoRadar LPC's latency is acceptable; on a
long-frame sensor like Kiem's it is the limiting factor.

% -------------------------------------------------------------------
\section{Throughput}
\label{sec:hw_throughput}

With a 128-bit input beat at clock $f$, the theoretical throughput is
$128 f / \mathrm{II}$. Table~\ref{tab:hw_throughput} sets it and the measured
rate against what the sensors require.

\begin{table}[htbp]
\centering
\caption[Throughput: measured, theoretical and required]{Throughput in decimal
Gbit/s. Measured figures include pipeline fill and drain: on the board for the
final designs, in RTL co-simulation for the replica. Kiem's
\SI{11.92}{Gibit/s} is \SI{12.8}{Gbit/s}, theoretical.}
\label{tab:hw_throughput}
\small
\begin{tabular}{@{}lrr@{}}
\toprule
& Theoretical & Measured\\
\midrule
DRHE, $\mathrm{II}=2$ (board) & 6.40 & \textbf{5.12}\\
LPC, frame average (board) & --- & \textbf{0.97}\\
Replica of Kiem, $\mathrm{II}=1$ (co-sim) & 12.80 & \HW{ablk}{thrmeas}\\
Kiem \cite{kiem2025} & 12.80 & not reported\\
\midrule
\emph{Required:} ColoRadar, 4 RX & \multicolumn{2}{l}{0.209}\\
\emph{Required:} ColoRadar, all 16 RX & \multicolumn{2}{l}{0.835}\\
\emph{Required:} Kiem's sensor & \multicolumn{2}{l}{3.20}\\
\bottomrule
\end{tabular}
\end{table}

Kiem's \SI{11.92}{Gibit/s} is stream width times clock, not a measurement.
Theoretical against theoretical, the gap is exactly the factor of two between
$\mathrm{II} = 2$ and $\mathrm{II} = 1$; the further \SI{20}{\percent} in the
measured figure is the per-ramp pipeline drain, which the replica shows too.
The comparison that matters is with the sensor: DRHE accepts data 24.5 times
faster than four ColoRadar channels produce it and 1.6 times faster than Kiem's
sensor, with no buffer. That headroom could let one instance serve all sixteen
channels, or run at a slower clock. LPC's \SI{0.97}{Gbit/s} is a frame
average; its ingest pass alone accepts \SI{12.8}{Gbit/s}, and only for a sensor
delivering frames back to back would its 1.16 times headroom over ColoRadar's
16-channel rate matter. At $\mathrm{II} = 2$ DRHE accepts data at half the rate
a streaming FFT core can produce it within a ramp, so the FFT must be allowed
to pause, or a FIFO of half a block RAM must sit between them.

% -------------------------------------------------------------------
\section{Resources against the device}
\label{sec:hw_resources}

\begin{table}[htbp]
\centering
\caption[Resources as a fraction of two devices]{Post-route resources of each
core as a percentage of the \texttt{XCKU5P} (\num{216960} LUT, \num{433920} FF,
480 BRAM, \num{1824} DSP) and of Kiem's \texttt{XCZU3CG} (\num{70560} LUT,
\num{141120} FF, 216 BRAM, 360 DSP). BRAM in 36\,Kb blocks. Final cores are the
in-system figures of Table~\ref{tab:resources}; the replica is out of context.}
\label{tab:hw_device}
\small
\setlength{\tabcolsep}{4pt}
\begin{tabular}{@{}lrrrrrrrr@{}}
\toprule
& \multicolumn{4}{c}{\% of \texttt{XCKU5P}} & \multicolumn{4}{c}{\% of \texttt{XCZU3CG}}\\
\cmidrule(lr){2-5}\cmidrule(lr){6-9}
Core & LUT & FF & BRAM & DSP & LUT & FF & BRAM & DSP\\
\midrule
DRHE compressor   & 16 & 5 & 8 & 7 & 50 & 17 & 19 & 37\\
DRHE decompressor & 7 & 2 & 7 & 3 & 20 & 7 & 15 & 13\\
LPC compressor    & 16 & 5 & 22 & 9 & 48 & 14 & 50 & 44\\
LPC decompressor  & 7 & 2 & 0 & 1.3 & 22 & 5 & 0 & 6\\
Replica of Kiem   & 8 & 1 & 2 & 3 & 24 & 4 & 5 & 13\\
Kiem's IP \cite{kiem2025} & 21 & 2 & 2 & 2 & 65 & 7 & 5 & 12\\
\bottomrule
\end{tabular}
\end{table}

Each compressor uses about a sixth of the \texttt{XCKU5P}'s logic, and encoder
and decoder together under a quarter (Table~\ref{tab:hw_device}). Every
compressor also fits Kiem's \texttt{XCZU3CG}, and the DRHE compressor uses
fewer of its LUTs than his own IP does (50 against \SI{65}{\percent}). His
complete system used \num{59353} LUTs, \num{13461} of them outside his
compressor; with DRHE's compressor in its place the system would use about
\SI{69}{\percent} of the device, inside the usual 70--\SI{80}{\percent}
congestion guideline. The decompressors sit at the central processor, not
beside the sensor, so they are not part of this budget.

The cores compress four of ColoRadar's sixteen receive channels. Four DRHE
compressors would fit the \texttt{XCKU5P} (\SI{65}{\percent} of its LUTs) but
not the \texttt{XCZU3CG}. One time-multiplexed instance, using its throughput
headroom, keeps the logic of one and needs 160 blocks for the state, which
fits Kiem's device only tightly. LPC's memory is its limit: at about 100 blocks
per four channels, sixteen channels need over \SI{80}{\percent} of the
\texttt{XCKU5P}'s block RAM. The largest
single saving available is DRHE's arithmetic: Kiem's fixed-point types would
roughly halve its logic, cut its DSPs to a third and its block RAM to half, for
\SI{0.07}{\percent} of compression ratio (Section~\ref{sec:hw_ladder}).

% -------------------------------------------------------------------
\section{Every difference from Kiem, accounted for}
\label{sec:hw_ledger}

Table~\ref{tab:hw_ledger} lists every number in which this thesis's DRHE and
Kiem's differ, whether the published numbers measure the same thing, and the
cause.

\begin{table}[htbp]
\centering
\caption[Every difference from Kiem's design]{Every difference between this
thesis's DRHE and Kiem's published design. ``Final'' is the lag-$\nTx$
compressor on the board; the replica is post-route, out of context. ``Valid''
says whether the published numbers measure the same thing.}
\label{tab:hw_ledger}
\footnotesize
\renewcommand{\arraystretch}{1.2}
\setlength{\tabcolsep}{3pt}
\begin{tabular}{@{}>{\raggedright}p{0.11\textwidth}>{\raggedleft}p{0.1\textwidth}>{\raggedleft}p{0.09\textwidth}>{\raggedleft}p{0.09\textwidth}>{\raggedright}p{0.07\textwidth}>{\raggedright\arraybackslash}p{0.43\textwidth}@{}}
\toprule
Metric & Final & Replica & Kiem & Valid & Cause of the difference\\
\midrule
LUT & \num{35444} & \HW{ablk}{LUT} & \num{45892} & yes & Against the replica: floating-point operators, then the packer. Replica against Kiem: about \SI{5}{\percent} explained.\\
FF & \num{23528} & \HW{ablk}{FF} & \num{9243} & yes & Registers inside float operators and along a 58-stage pipeline. Replica against Kiem: about \SI{45}{\percent} explained.\\
DSP & 132 & \HW{ablk}{DSP} & 44 & yes & Every DSP is a float operator; sharing at $\mathrm{II}=2$ halves the count.\\
BRAM & 40 & \HW{ablk}{BRAMT} & 10 & yes & 32-bit state words need twice the memory of 16-bit ones; per-transmitter state doubles it again.\\
II & 2 & 1 & 1 & yes & One extra write per cycle to clear the state.\\
Depth & 58 & \HW{ablk}{hDepth} & 23 & yes & Closed loop (26 cycles), then float operator latency.\\
$F_{\max}$ & 100 met & \HW{ablk}{fmax} & 100 met & partly & All meet \SI{100}{\mega\hertz}; Kiem on speed grade $-1$.\\
Throughput & 5.12 meas. & \HW{ablk}{thrmeas} & 12.8 theor. & theor.\ only & $\mathrm{II}=2$ against $\mathrm{II}=1$.\\
Added latency & \SI{13.1}{\micro\second}; \SI{21.3}{\micro\second} with FFT & --- & \SI{79.4}{\micro\second} & with FFT & Same definition; his chain, like the \SI{21.3}{\micro\second} figure, includes the range FFT.\\
Compression ratio & 3.75, 4.22 & --- & 3.17 & no & Different sensor and data scaling (Chapter~\ref{ch:vskiem}).\\
\bottomrule
\end{tabular}
\end{table}

Measured post-route, as Kiem's figures are, this thesis's DRHE compressor uses
fewer LUTs than his IP, meets the same clock and fits his device at half its
logic. Against his design rebuilt in the same tool it is larger, by exactly
what its distinctive decisions cost: floating-point arithmetic (most of the
logic, every extra DSP, twice the block RAM), a less economical packer (about
\num{4200} LUTs), a closed loop (26 cycles in floating point, almost nothing in
fixed point), one extra state write ($\mathrm{II} = 2$, half the throughput and
half the arithmetic), and per-transmitter state (twice the block RAM again).
The closed loop guarantees the encoder and decoder never diverge, the
per-transmitter state is what makes prediction pay on a TDM-MIMO sensor, and
the floating-point datapath buys only \SI{0.07}{\percent} of ratio.
